% =====================================================================
%  FOCUS! - Tutorial: abrir el portal web y revisar todo en Azure
%  Compilar en Overleaf con pdfLaTeX (Menu > Compiler > pdfLaTeX).
% =====================================================================
\documentclass[11pt,a4paper]{article}

\usepackage[utf8]{inputenc}
\usepackage[T1]{fontenc}
\usepackage[spanish,es-noshorthands]{babel}
\usepackage{lmodern}
\usepackage[margin=2.3cm]{geometry}
\usepackage{xcolor}
\usepackage{hyperref}
\usepackage{listings}
\usepackage{booktabs}
\usepackage{tabularx}
\usepackage{enumitem}
\usepackage{fancyhdr}
\usepackage{titlesec}
\usepackage[most]{tcolorbox}
\usepackage{tikz}

% ---------- Colores de la app (res/values/colors.xml) ----------
\definecolor{focusnavy}{HTML}{0F172A}
\definecolor{focusorange}{HTML}{F59E0B}
\definecolor{focusslate}{HTML}{64748B}
\definecolor{focusgreen}{HTML}{10B981}
\definecolor{focusbg}{HTML}{F8FAFC}

% ---------- Logo de Focus! (ic_logo.xml + ic_orbit.xml), igual que en los documentos anteriores ----------
\newcommand{\logofocus}[1][1]{%
  \begin{tikzpicture}[x=#1*0.035cm, y=#1*0.035cm, baseline=(c.base)]
    \coordinate (c) at (0,0);
    \draw[focusorange!25, line width=#1*0.6pt] (0,0) circle (64);   % órbita
    \fill[focusorange] (41,52) circle (3);                          % punto superior
    \fill[focusorange] (-36,-58) circle (2.5);                      % punto inferior
    \fill[focusnavy] (0,0) circle (38.5);                           % anillo
    \fill[white]     (0,0) circle (20.125);                         % núcleo
    \fill[focusnavy] (0,0) circle (5.25);                           % centro
  \end{tikzpicture}%
}
\newcommand{\marcafocus}[1]{{\sffamily\bfseries\fontsize{#1}{#1}\selectfont
  \color{focusnavy}FOCUS\textcolor{focusorange}{!}}}

% ---------- Colores del portal web ----------
\definecolor{ambar}{HTML}{F59E0B}
\definecolor{ambaroscuro}{HTML}{B45309}
\definecolor{marino}{HTML}{131C2E}
\definecolor{gris}{HTML}{64748B}
\definecolor{fondo}{HTML}{F8FAFC}
\definecolor{verde}{HTML}{047857}
\definecolor{rojo}{HTML}{B91C1C}

\hypersetup{colorlinks=true, linkcolor=ambaroscuro, urlcolor=ambaroscuro, pdftitle={FOCUS! - Tutorial del portal y Azure}}

\titleformat{\section}{\Large\bfseries\color{marino}}{\thesection}{0.6em}{}
\titleformat{\subsection}{\large\bfseries\color{ambaroscuro}}{\thesubsection}{0.6em}{}

\pagestyle{fancy}
\fancyhf{}
\lhead{\small\color{gris}FOCUS! \textperiodcentered{} Tutorial del portal y de Azure}
\rhead{\small\color{gris}\thepage}
\renewcommand{\headrulewidth}{0.4pt}

% ---------- Código (acentos permitidos dentro de los listados) ----------
\lstset{
  basicstyle=\ttfamily\small, backgroundcolor=\color{fondo}, frame=single, rulecolor=\color{gris!40},
  breaklines=true, columns=fullflexible, keepspaces=true, showstringspaces=false,
  keywordstyle=\color{ambaroscuro}\bfseries, commentstyle=\color{gris}\itshape, stringstyle=\color{verde},
  literate={á}{{\'a}}1 {é}{{\'e}}1 {í}{{\'i}}1 {ó}{{\'o}}1 {ú}{{\'u}}1 {ñ}{{\~n}}1
           {Á}{{\'A}}1 {É}{{\'E}}1 {Í}{{\'I}}1 {Ó}{{\'O}}1 {Ú}{{\'U}}1 {Ñ}{{\~N}}1
           {¿}{{?`}}1 {¡}{{!`}}1 {·}{{\textperiodcentered}}1 {—}{{---}}1 {–}{{--}}1 {›}{{>}}1
}
\lstdefinelanguage{texto}{}

% ---------- Cajas ----------
\newtcolorbox{nota}[1][Nota]{colback=ambar!8, colframe=ambar, fonttitle=\bfseries, title=#1, boxrule=0.6pt, arc=2mm}
\newtcolorbox{alerta}[1][Importante]{colback=rojo!6, colframe=rojo, fonttitle=\bfseries, title=#1, boxrule=0.6pt, arc=2mm}
\newtcolorbox{paso}[1]{colback=white, colframe=marino, fonttitle=\bfseries, title=#1, boxrule=0.6pt, arc=2mm}

\newcommand{\ruta}[1]{\texttt{\color{marino}#1}}
\setlength{\emergencystretch}{3em}  % evita líneas que se salen del margen
\newcommand{\menu}[1]{\textsf{\textbf{#1}}}

\begin{document}

% =====================================================================
% PORTADA (la misma de los documentos anteriores de Focus!)
\begin{titlepage}
  \begin{tikzpicture}[remember picture, overlay]
    \fill[focusbg] (current page.south west) rectangle (current page.north east);
    \fill[focusorange!10] ([xshift=-1.5cm,yshift=-2.5cm]current page.north east) circle (4.5cm);
    \fill[focusgreen!7] ([xshift=2cm,yshift=2.5cm]current page.south west) circle (5.5cm);
  \end{tikzpicture}
  \centering
  {\large\sffamily\color{focusnavy}Universidad Autónoma de Querétaro}\\[0.2cm]
  {\sffamily\color{focusslate}Facultad de Informática \textperiodcentered{} Ingeniería de Software}\\[1.8cm]
  \logofocus[1]\\[0.5cm]
  \marcafocus{40}\\[0.3cm]
  {\large\color{focusslate}Menos pantalla, \textcolor{focusnavy}{\textbf{más enfoque.}}}\\[1.8cm]
  {\sffamily\bfseries\LARGE\color{focusnavy}Tutorial del portal web}\\[0.15cm]
  {\sffamily\bfseries\LARGE\color{focusnavy}y de Microsoft Azure}\\[0.35cm]
  {\large\sffamily\color{focusslate}Abrir la web, usar la API y revisar cada servicio}\\[2cm]
  \begin{tabular}{rl}
    \textbf{Alumno:} & Luis Angel Barrera Barrera \\
    \textbf{Profesor:} & Gonzalez Gutierrez Fidel \\
    \textbf{Grupo:} & 31 \textperiodcentered{} Clave 2075 \\
    \textbf{Portal web:} & \url{https://green-moss-07816e210.4.azurestaticapps.net} \\
    \textbf{Repositorio:} & \url{https://github.com/lbarrera27-pixel/focus_web} \\
  \end{tabular}\\[1.5cm]
  {\sffamily\color{focusslate}Noviembre de 2026}
\end{titlepage}

\tableofcontents
\newpage

% =====================================================================
\section{Cómo está armado el proyecto}

FOCUS! ayuda a reducir el tiempo de pantalla. La app Android mide el uso del teléfono y muestra pausas
antes de abrir apps; el \textbf{portal web} muestra el plan, el riesgo de recaída, los reportes y los
ajustes de la cuenta. Todo vive en Microsoft Azure, dentro del grupo de recursos \texttt{rg-focus}.

\begin{center}
\renewcommand{\arraystretch}{1.25}
\begin{tabularx}{\textwidth}{>{\bfseries}l l X}
\toprule
Parte & Servicio de Azure & Nombre del recurso \\
\midrule
Portal web (frontend) & Static Web Apps & \texttt{focus-web-luisbb27} \\
API (backend FastAPI) & App Service (Linux, Python 3.12) & \texttt{focus-api-luisbb27} (plan \texttt{plan-focus}, B1) \\
Base de datos & Azure SQL Database (serverless) & servidor \texttt{focus-sql-luisbb27}, base \texttt{focusdb} \\
Fotos de perfil & Blob Storage & cuenta \texttt{focusstluisbb27}, contenedor \texttt{fotos} \\
Función serverless & Azure Functions & \texttt{focus-func-luisbb27} (\ruta{/api/ajuste-semanal}) \\
IA conversacional & Claude (Anthropic) & la usa la API para el Compañero Foqui \\
\bottomrule
\end{tabularx}
\end{center}

\noindent El flujo es: \textbf{navegador o teléfono} $\rightarrow$ \textbf{API} (con el header
\texttt{Authorization: Bearer}) $\rightarrow$ \textbf{Azure SQL} y \textbf{Blob Storage}. La API también llama a
Claude para las respuestas de Foqui.

\begin{nota}[La base de datos se duerme]
Azure SQL está en modo \emph{serverless}: se pausa tras una hora sin uso para ahorrar créditos.
La primera petición después de una pausa tarda unos 30--60 segundos; mientras despierta, la API
responde \texttt{503} y el portal reintenta solo. Si algo tarda al principio, espera y recarga.
\end{nota}

% =====================================================================
\section{Abrir el portal web}

\begin{paso}{Paso 1 \textperiodcentered{} Entrar}
\begin{enumerate}[leftmargin=*]
  \item Abre \url{https://green-moss-07816e210.4.azurestaticapps.net} en Chrome, Edge o Firefox.
  \item Escribe el \textbf{correo} y la \textbf{contraseña} de la cuenta. Para la presentación se usa la cuenta
        con datos de ejemplo \texttt{demo2@focus.app} (la contraseña la tiene el equipo).
  \item Marca \menu{Mantener la sesión iniciada} si no quieres volver a escribirla al cerrar la pestaña.
  \item Pulsa \menu{Iniciar sesión}.
\end{enumerate}
\end{paso}

\begin{nota}[Sin cuenta]
El botón \menu{Ver demo sin cuenta} abre el portal con datos de demostración que viven solo en el
navegador. Sirve para enseñar las pantallas, pero \textbf{no} lee ni escribe nada en Azure.
\end{nota}

\subsection{Qué hay en cada sección}

\begin{center}
\renewcommand{\arraystretch}{1.2}
\begin{tabularx}{\textwidth}{>{\bfseries}l X l}
\toprule
Sección & Qué muestra & Rutas de la API \\
\midrule
Inicio & Panel general: uso por día, apps, pausas, plan y dispositivo & \ruta{/estadisticas}, \ruta{/eventos} \\
Mi plan & Límite diario, semana actual, próximo ajuste y simulador & \ruta{/planes}, \ruta{/planes/\{n\}/ajustes} \\
Riesgo de recaída & Probabilidad de pasarse del límite la próxima semana & \ruta{/riesgo} \\
Uso por app & Tiempo por app, categorías y mapa de calor día $\times$ hora & \ruta{/eventos}, \ruta{/apps} \\
Apps y sitios & Apps restringidas, sitios bloqueados, límites y reglas & \ruta{/apps}, \ruta{/dominios} \\
Foqui & Chat con el compañero de IA sobre tu semana & \ruta{/companero/mensaje} \\
Reportes & Resúmenes semanales, PDF y exportar CSV & \ruta{/eventos}, \ruta{/planes} \\
Ajustes & Cuenta, notificaciones, privacidad y dispositivo & \ruta{/usuarios/me}, \ruta{/configuracion/me} \\
\bottomrule
\end{tabularx}
\end{center}

\noindent Al pulsar tu nombre (abajo a la izquierda) se abre el \textbf{menú de la cuenta}: Mi perfil,
Ajustes y reglas, Exportar mis datos, Ayuda y soporte y Cerrar sesión. Una dirección que no existe
(por ejemplo \texttt{\#/abc}) muestra la página \emph{No encontramos esta página}.

\subsection{Recorrido sugerido para la presentación}
\begin{enumerate}[leftmargin=*]
  \item \textbf{Uso por app}: el mapa de calor y la franja más crítica (más uso de noche).
  \item \textbf{Mi plan}: el límite baja cada semana según las reglas difusas de la API.
  \item \textbf{Riesgo de recaída}: la regresión logística y los factores que la explican.
  \item \textbf{Foqui}: escribe \emph{``Resúmeme mi semana''}; la respuesta la genera Claude con tus números.
  \item \textbf{Reportes}: descarga el PDF de una semana y exporta el CSV.
  \item \textbf{Ajustes \textrightarrow{} Dispositivo}: el teléfono vinculado, su batería y la actividad de hoy.
  \item \textbf{Ajustes \textrightarrow{} Cuenta}: cambia la foto (se guarda en Blob Storage).
\end{enumerate}

% =====================================================================
\section{Usar la API}

\subsection{Swagger: probar desde el navegador}
\begin{enumerate}[leftmargin=*]
  \item Abre \url{https://focus-api-luisbb27.azurewebsites.net/docs}.
  \item Prueba primero \ruta{GET /salud} con \menu{Try it out} \textrightarrow{} \menu{Execute}. Debe responder:
\begin{lstlisting}
{"status":"ok","base_datos":"ok","motor":"mssql"}
\end{lstlisting}
  \item Para las rutas privadas entra con \ruta{POST /auth/login} y copia el \texttt{access\_token} de la respuesta.
  \item Clic en \menu{Authorize}, arriba a la derecha, pega el token y confirma. Swagger lo manda en cada petición.
\end{enumerate}

\begin{nota}[Los números de las rutas son por usuario]
\ruta{GET /pausas/7} con el token del usuario 8 es la 7.\textordmasculine{} pausa de ese usuario, no la pausa
con id 7 de la tabla. Sin token la API responde \texttt{422}; con un token inválido o vencido, o de un
usuario desactivado, \texttt{401}.
\end{nota}

\subsection{Desde PowerShell}
\begin{lstlisting}[language=texto]
$api = "https://focus-api-luisbb27.azurewebsites.net"
$cuerpo = '{"email":"demo2@focus.app","password":"TU_CONTRASENA"}'
$login  = Invoke-RestMethod "$api/auth/login" -Method Post -ContentType "application/json" -Body $cuerpo
$h   = @{ Authorization = "Bearer $($login.access_token)" }

Invoke-RestMethod "$api/salud"
Invoke-RestMethod "$api/apps" -Headers $h
Invoke-RestMethod "$api/planes/activo" -Headers $h
Invoke-RestMethod "$api/riesgo" -Headers $h
Invoke-RestMethod "$api/tokens-dispositivo" -Headers $h
Invoke-RestMethod "$api/estadisticas?desde=2026-10-28&hasta=2026-11-03" -Headers $h
\end{lstlisting}

\subsection{Con Postman}
Postman tiene su propia sección, paso a paso: la sección \ref{sec:postman}.

% =====================================================================
\section{Probar las APIs con Postman}\label{sec:postman}

\begin{nota}[Postman no ``levanta'' la API]
Postman solo \textbf{hace peticiones}; no enciende el servidor. La API ya está encendida en Azure
(\url{https://focus-api-luisbb27.azurewebsites.net}), así que normalmente solo hay que apuntar Postman ahí.
Si quieres probar la API en tu computadora, primero la enciendes con los comandos de la sección
\ref{sec:local} y después apuntas Postman a \texttt{http://127.0.0.1:8000}.
\end{nota}

\subsection{Las dos colecciones}
\begin{center}
\renewcommand{\arraystretch}{1.25}
\begin{tabularx}{\textwidth}{>{\ttfamily\raggedright\arraybackslash}p{5.2cm} X}
\toprule
\normalfont\bfseries Archivo & \textbf{Qué trae} \\
\midrule
Focus\_API\_v2\newline .postman\_collection.json & Todas las rutas de la API (más de 70 operaciones), la función serverless, las técnicas de IA y los casos de error. \\
Focus\_API\_nuevas\newline .postman\_collection.json & Solo las rutas nuevas y las que usa el portal nuevo: 25 peticiones y 33 pruebas. \\
\bottomrule
\end{tabularx}
\end{center}
Las dos están en la carpeta \texttt{postman/} del repositorio y en el ZIP de la API.

\subsection{Paso a paso}
\begin{paso}{1 \textperiodcentered{} Instalar e importar}
\begin{enumerate}[leftmargin=*]
  \item Instala Postman de escritorio desde \url{https://www.postman.com/downloads/} (la versión web
        no puede llamar a \texttt{127.0.0.1} sin instalar el \emph{Postman Agent}).
  \item \menu{Import} (arriba a la izquierda) \textrightarrow{} arrastra el archivo \texttt{.json} de la colección.
\end{enumerate}
\end{paso}

\begin{paso}{2 \textperiodcentered{} Elegir a qué API apuntar}
\begin{enumerate}[leftmargin=*]
  \item Abre la colección \textrightarrow{} pestaña \menu{Variables}.
  \item En \texttt{baseUrl} pon una de estas dos y guarda con \menu{Ctrl+S}:
    \begin{itemize}
      \item Azure: \texttt{https://focus-api-luisbb27.azurewebsites.net}
      \item Tu computadora: \texttt{http://127.0.0.1:8000}
    \end{itemize}
  \item En la colección completa, \texttt{funcUrl} es la función serverless:\\
        \texttt{https://focus-func-luisbb27.azurewebsites.net}.
\end{enumerate}
\end{paso}

\begin{paso}{3 \textperiodcentered{} Correr todo de una vez (Runner)}
\begin{enumerate}[leftmargin=*]
  \item Clic derecho en la colección \textrightarrow{} \menu{Run collection} \textrightarrow{} \menu{Run}.
  \item La colección hace sola el registro de un usuario de prueba, inicia sesión, guarda el token en la
        variable \texttt{token} y agrega el header \texttt{Authorization} a cada petición.
  \item Si la base de Azure estaba dormida, \texttt{GET /salud} espera 30 segundos y reintenta solo.
  \item Al final, la carpeta \emph{Limpieza} borra lo que se creó y desactiva el usuario de prueba.
  \item Resultado esperado de la colección de rutas nuevas: \textbf{25 peticiones, 33 pruebas en verde} y una
        \emph{omitida} (la de la foto, que necesita que elijas una imagen).
\end{enumerate}
\end{paso}

\begin{paso}{4 \textperiodcentered{} Mandar una petición a mano}
\begin{enumerate}[leftmargin=*]
  \item Abre \texttt{GET /salud} \textrightarrow{} \menu{Send}. Debe responder \texttt{\{"status":"ok","base\_datos":"ok"\}}.
  \item Para ver los datos de una cuenta que ya existe, entra con \texttt{POST /auth/login}
        (Body \textrightarrow{} raw \textrightarrow{} JSON con \texttt{email} y \texttt{password}). La respuesta trae \texttt{access\_token}.
  \item Copia ese valor en la variable \texttt{token} de la colección. Desde ahí todas las
        peticiones son de ese usuario. El token dura 30 días; cuando venza, vuelve a entrar.
  \item Para subir una foto: \texttt{POST /usuarios/me/foto/archivo} \textrightarrow{} \menu{Body}
        \textrightarrow{} \menu{form-data} \textrightarrow{} en la llave \texttt{archivo} elige \menu{File} y una imagen
        JPG, PNG o WEBP de hasta 5 MB \textrightarrow{} \menu{Send}. Se guarda en Blob Storage.
\end{enumerate}
\end{paso}

\subsection{Las rutas nuevas}
\begin{center}
\small
\renewcommand{\arraystretch}{1.25}
\begin{tabularx}{\textwidth}{l >{\ttfamily\raggedright\arraybackslash}p{5.7cm} X}
\toprule
\textbf{Método} & \normalfont\bfseries Ruta & \textbf{Para qué} \\
\midrule
PUT & /tokens-dispositivo/\{n\}/estado & El teléfono reporta modelo, versión y batería. Cuerpo: \texttt{\{"modelo":"Pixel 7 Pro", "version\_app":"1.8", "bateria":78, "cargando":false\}}. Lo que no se manda se conserva. \\
GET & /tokens-dispositivo/\{n\}/estado & Último estado que reportó (404 si todavía no reporta). \\
GET & /tokens-dispositivo & Ahora incluye \texttt{estado} en cada teléfono. \\
POST & /dominios/\{n\}/intentos & El teléfono avisa que detuvo un sitio bloqueado. Cuerpo opcional: \texttt{\{"fecha": "...UTC"\}}. \\
GET & /dominios/intentos?dias=7 & Intentos por sitio y por día (de 1 a 90 días). \\
POST & /usuarios/me/foto/archivo & Sube la foto de perfil a Blob Storage (form-data, campo \texttt{archivo}). \\
GET / PUT & /configuracion/me & Ajustes del portal (preferencias, avisos, privacidad) como JSON. \\
POST & /companero/mensaje & Foqui responde con IA; \texttt{fuente} dice si usó Claude (\texttt{llm}) o el respaldo. \\
GET & /riesgo & Probabilidad de pasarse del límite la próxima semana (404 sin plan con ajustes). \\
\bottomrule
\end{tabularx}
\end{center}

\subsection{Desde la terminal (newman)}
Corre la misma colección sin abrir Postman (necesita Node.js):
\begin{lstlisting}[language=texto]
# Contra Azure (usa la baseUrl guardada en la colección)
npx newman@6.2.1 run postman/Focus_API_nuevas.postman_collection.json

# Contra tu computadora
npx newman@6.2.1 run postman/Focus_API_nuevas.postman_collection.json --env-var baseUrl=http://127.0.0.1:8000
\end{lstlisting}

\subsection{Encender la API en tu computadora}\label{sec:local}
Desde la carpeta del repositorio (o del ZIP descomprimido), en PowerShell:
\begin{lstlisting}[language=texto]
python -m venv .venv
.venv\Scripts\activate
pip install -r requirements.txt
uvicorn main:app --reload          # queda escuchando en http://127.0.0.1:8000
\end{lstlisting}
Sin configurar nada usa una base SQLite local (\texttt{focus.db}); Swagger queda en
\url{http://127.0.0.1:8000/docs}. Deja esa ventana abierta mientras pruebas y ciérrala con \menu{Ctrl+C}.

\subsection{Si algo sale mal}
\begin{center}
\renewcommand{\arraystretch}{1.25}
\begin{tabularx}{\textwidth}{>{\raggedright\arraybackslash}p{4.6cm} X}
\toprule
Respuesta & Qué significa \\
\midrule
\texttt{503} & La base de Azure está despertando. Espera 30 s y vuelve a enviar. \\
\texttt{401} & El token es inválido o venció, o el usuario está desactivado. Vuelve a entrar con \texttt{POST /auth/login} y actualiza la variable \texttt{token}. \\
\texttt{422} & Falta un dato o tiene un formato inválido (por ejemplo, batería mayor a 100). \\
\texttt{404} & Ese número no existe para ese usuario (los números son por usuario). \\
\emph{Could not send request} & Con \texttt{127.0.0.1} la API local no está encendida; con Azure, revisa tu internet. \\
\bottomrule
\end{tabularx}
\end{center}

% =====================================================================
\section{Revisar todo en el portal de Azure}

Entra a \url{https://portal.azure.com} con la cuenta de \emph{Azure for Students} y abre
\menu{Grupos de recursos} \textrightarrow{} \texttt{rg-focus}. Ahí están todos los recursos del proyecto.

\subsection{API: App Service \texttt{focus-api-luisbb27}}
\begin{center}
\renewcommand{\arraystretch}{1.2}
\begin{tabularx}{\textwidth}{>{\bfseries}l X}
\toprule
Dónde & Qué revisar \\
\midrule
Información general & Estado \emph{En ejecución}, la URL y los botones \menu{Reiniciar} / \menu{Detener}. \\
Supervisión \textrightarrow{} Secuencia de registro & Los mensajes de la API en vivo (peticiones y errores). \\
Configuración \textrightarrow{} Variables de entorno & \texttt{AZURE\_SQL\_*}, \texttt{AZURE\_STORAGE\_*} y \texttt{ANTHROPIC\_API\_KEY}. Los valores son secretos: no los compartas. \\
Centro de implementación \textrightarrow{} Registros & Cada despliegue con su fecha y si terminó bien. \\
Supervisión \textrightarrow{} Métricas & Peticiones, errores HTTP 5xx y tiempo de respuesta. \\
\bottomrule
\end{tabularx}
\end{center}

\subsection{Base de datos: Azure SQL \texttt{focusdb}}
\begin{paso}{Abrir el editor de consultas}
\begin{enumerate}[leftmargin=*]
  \item En \texttt{rg-focus} abre la base \texttt{focusdb} (servidor \texttt{focus-sql-luisbb27}).
  \item En \menu{Información general} revisa el estado: \emph{En línea} o \emph{En pausa}. Si está en pausa,
        se despierta sola al abrir el editor (tarda un minuto).
  \item Abre \menu{Editor de consultas (versión preliminar)}.
  \item Inicia sesión con \textbf{Autenticación de SQL Server}: usuario \texttt{focusadmin} y la contraseña
        que generó el script de despliegue (está en \texttt{deploy\textbackslash salida.env}, que no se sube a GitHub).
  \item Si aparece un aviso de firewall, pulsa \menu{Permitir la dirección IP del cliente} y vuelve a entrar.
\end{enumerate}
\end{paso}

\noindent Las 13 tablas del proyecto:
\begin{center}
\small
\begin{tabularx}{\textwidth}{>{\ttfamily}l X}
\toprule
usuario & Cuentas (correo, nombre, contraseña cifrada, foto, activo) \\
app & Apps que mide el teléfono y su límite de aperturas \\
evento\_uso & Cada uso de una app (inicio, fin, duración) \\
franja & Catálogo de franjas horarias \\
plan / ajuste\_plan & Plan del usuario y su ajuste de cada semana (reglas difusas) \\
pausa & Pausas de fricción mostradas y si se ignoraron \\
dominio\_bloqueado / intento\_sitio & Sitios bloqueados y cada intento de abrirlos \\
sesion\_foco & Sesiones de Modo Foco \\
token\_dispositivo / estado\_dispositivo & Teléfonos vinculados y su último estado (modelo, versión, batería) \\
configuracion\_usuario & Ajustes del portal (JSON por usuario) \\
\bottomrule
\end{tabularx}
\end{center}

\noindent Consultas útiles para mostrar los datos:
\begin{lstlisting}[language=SQL]
-- Tablas y cuántos registros tiene cada una
SELECT t.name AS tabla, SUM(p.rows) AS filas
FROM sys.tables t JOIN sys.partitions p ON p.object_id = t.object_id AND p.index_id IN (0, 1)
GROUP BY t.name ORDER BY t.name;

-- Usuarios registrados
SELECT id, email, nombre, activo, fecha_registro FROM usuario ORDER BY id;

-- Uso de la cuenta de ejemplo (usuario 8) por app, en minutos
SELECT a.nombre, COUNT(*) AS aperturas, SUM(e.duracion_seg) / 60 AS minutos
FROM evento_uso e JOIN app a ON a.id = e.app_id
WHERE e.usuario_id = 8
GROUP BY a.nombre ORDER BY minutos DESC;

-- Plan activo y sus ajustes semanales
SELECT p.numero, p.fecha_inicio, p.limite_base, j.semana, j.uso_meta, j.adherencia, j.cambio_pct
FROM plan p LEFT JOIN ajuste_plan j ON j.plan_id = p.id
WHERE p.usuario_id = 8 AND p.activo = 1 ORDER BY j.semana;

-- Pausas de hoy (la fecha se guarda en UTC)
SELECT numero, tipo, duracion_seg, ignorada, fecha FROM pausa
WHERE usuario_id = 8 AND CAST(fecha AS date) = CAST(SYSUTCDATETIME() AS date);

-- Teléfono vinculado y lo último que reportó
SELECT t.numero, t.plataforma, t.activo, s.modelo, s.version_app, s.bateria, s.fecha
FROM token_dispositivo t LEFT JOIN estado_dispositivo s ON s.token_id = t.id
WHERE t.usuario_id = 8;
\end{lstlisting}

\begin{alerta}[Solo consultar]
Usa \texttt{SELECT} para revisar. No borres ni cambies filas desde el editor: la API lleva la
numeración por usuario y la integridad de los datos. Para cambiar algo, usa el portal o la API.
\end{alerta}

\subsection{Fotos: Blob Storage \texttt{focusstluisbb27}}
\menu{Almacenamiento de datos} \textrightarrow{} \menu{Contenedores} \textrightarrow{} \texttt{fotos}. Cada
usuario tiene su carpeta (por ejemplo \texttt{usuario-8/}). Al cambiar la foto en
\menu{Ajustes \textrightarrow{} Cuenta} aparece un archivo nuevo y la API guarda su enlace en \texttt{usuario.foto\_url}.

\subsection{Función serverless: \texttt{focus-func-luisbb27}}
\menu{Funciones} \textrightarrow{} \texttt{ajuste\_semanal} \textrightarrow{} \menu{Código y prueba}
permite ejecutarla desde el portal. Calcula el ajuste semanal con las mismas reglas difusas de la API;
\menu{Supervisar} muestra cada ejecución.

\subsection{Portal web: Static Web App \texttt{focus-web-luisbb27}}
En \menu{Información general} está la URL pública; en \menu{Entornos} aparece la versión publicada
(\emph{production}) con la fecha del último despliegue.

\subsection{Créditos}
\menu{Suscripciones} \textrightarrow{} \emph{Azure for Students} \textrightarrow{} \menu{Análisis de costos}
muestra cuánto gasta cada recurso. Lo que más consume es el plan B1 del App Service.

% =====================================================================
\section{Revisar desde la terminal (Azure CLI)}
Requiere \texttt{az login}. Los comandos son de solo lectura salvo el reinicio.
\begin{lstlisting}[language=texto]
az resource list -g rg-focus -o table                              # todos los recursos
az webapp show -g rg-focus -n focus-api-luisbb27 --query state       # ¿está corriendo?
az webapp log tail -g rg-focus -n focus-api-luisbb27                # registro en vivo (Ctrl+C para salir)
az webapp config appsettings list -g rg-focus -n focus-api-luisbb27 --query "[].name" -o tsv
az sql db show -g rg-focus -s focus-sql-luisbb27 -n focusdb --query "{estado:status, sku:currentSku.name}"
az storage blob list --account-name focusstluisbb27 -c fotos --auth-mode login -o table
az staticwebapp show -g rg-focus -n focus-web-luisbb27 --query defaultHostname
az webapp restart -g rg-focus -n focus-api-luisbb27                 # reiniciar la API
\end{lstlisting}

% =====================================================================
\section{Publicar cambios}
Desde PowerShell, en la carpeta del repositorio:
\begin{lstlisting}[language=texto]
.\deploy\02_desplegar_api.ps1     # sube la API, la reinicia y prueba /salud
.\deploy\05_desplegar_web.ps1     # compila el portal contra la API y lo publica
python deploy/06_datos_demo.py --api https://focus-api-luisbb27.azurewebsites.net --email demo2@focus.app --refrescar
\end{lstlisting}
El último comando pone al día la cuenta de ejemplo (eventos, pausas de hoy y semanas del plan).
Conviene correrlo el día de la presentación.

% =====================================================================
\section{Avance del proyecto}

\subsection{Lo que ya está hecho y publicado}
\begin{itemize}[leftmargin=*]
  \item \textbf{Portal web con el diseño de Figma (pantallas e1--e47)}: acceso, Mi plan, Riesgo de recaída,
        Uso por app, Apps y sitios, Foqui, Reportes, Ajustes de la cuenta y Cuenta (menú, perfil,
        desactivar, cerrar sesión, sesión expirada y página no encontrada).
  \item \textbf{API} con 13 tablas en capas (controladores, servicios, repositorios, modelos y esquemas),
        más de 70 operaciones, 94 pruebas automáticas y colección de Postman.
  \item \textbf{IA}: lógica difusa (ajuste semanal), regresión logística (riesgo), aprendizaje supervisado
        (pausas ignoradas), red neuronal (pronóstico de uso), algoritmo genético (horario de pausas),
        búsqueda A* (ruta del plan) y el Compañero Foqui con Claude.
  \item \textbf{Cambios recientes en la API} (todos solo agregan, no rompen nada de lo anterior):
        intentos de abrir sitios bloqueados (\ruta{/dominios/\{n\}/intentos}), estado del teléfono
        (\ruta{/tokens-dispositivo/\{n\}/estado}), subida de fotos desde el portal y la corrección
        para que los ajustes del plan se guarden en Azure SQL.
  \item \textbf{Portal}: preferencias de zona horaria, formato de hora (24/12 h) e inicio de semana (lunes o domingo),
        sin desplazamiento lateral desde 1024 px.
  \item \textbf{Datos de ejemplo} en la cuenta \texttt{demo2@focus.app} (usuario 8): 6 apps,
        unos 280 eventos, pausas, plan en la semana 4 con 3 ajustes, teléfono vinculado y 3 sitios bloqueados.
\end{itemize}

\subsection{Pendiente}
\begin{itemize}[leftmargin=*]
  \item \textbf{Cambiar la clave de Anthropic} (se mostró en un chat): crear una nueva en Claude Console,
        borrar la anterior y guardarla en la variable \texttt{ANTHROPIC\_API\_KEY} del App Service.
  \item \textbf{App Android} (repositorio \texttt{focus\_app}): enviar su modelo, versión y batería con
        \ruta{PUT /tokens-dispositivo/\{n\}/estado} para que las cuentas reales los muestren.
\end{itemize}

% =====================================================================
\section{Problemas comunes}
\begin{center}
\renewcommand{\arraystretch}{1.25}
\begin{tabularx}{\textwidth}{>{\raggedright\arraybackslash}p{5.2cm} X}
\toprule
Síntoma & Qué hacer \\
\midrule
El portal tarda o dice \emph{``La base de datos se está iniciando''} & La base serverless estaba en pausa. Espera 30--60 s y recarga. \\
Sale \emph{``Tu sesión terminó''} & La sesión expiró (30 min sin actividad si no marcaste ``Mantener la sesión'') o la cuenta se desactivó. Vuelve a entrar. \\
Foqui responde con mensajes genéricos & Falta o es inválida \texttt{ANTHROPIC\_API\_KEY} en el App Service. \\
Una cuenta se ve vacía & No tiene datos. Corre \texttt{06\_datos\_demo.py} o usa \texttt{demo2@focus.app}. \\
Después de desplegar la API sigue la versión anterior & \texttt{az webapp restart -g rg-focus -n focus-api-luisbb27}. \\
El editor de consultas no deja entrar & Agrega tu IP en el firewall del servidor SQL (\menu{Redes}). \\
\bottomrule
\end{tabularx}
\end{center}

\end{document}
